\section*{\fontsize{18pt}{18pt}\selectfont\MakeUppercase{\de}\\TÓM TẮT BÁO CÁO TUẦN}

Báo cáo này tổng hợp công việc nhóm thực hiện trong tuần \tuanbaocao{} của đề tài ``Xây dựng mạng Campus Network sử dụng SDN và SD-WAN''. Sau khi phần SDN đã có đủ số liệu ở tuần trước, trọng tâm tuần này là thu số liệu SD-WAN đầu tiên, kiểm thử tích hợp toàn hệ thống và xây dựng công cụ vận hành mạng tự động.

Kết quả chính của tuần:

\begin{itemize}
    \item \textbf{Đo kiểm SD-WAN:} failover link WAN tại Site 200 có RTO trung bình 7,96~s (7,68--8,25~s), PRR 76,75\%; so sánh hiệu năng bốn site cho thấy đường WAN qua SD-WAN không mất gói, delay 3,8--12~ms, jitter cao nhất ở Site 400.
    \item \textbf{Kiểm thử tích hợp toàn hệ thống (70 thiết bị):} SDN 6/6 OVS, firewall HA Active/Standby, SD-WAN 12/12 thiết bị và BFD 92/92, DHCP và truy cập Web/Mail/DNS từ cả ba chi nhánh, Internet 4/4 site -- đều ĐẠT.
    \item \textbf{Sửa hai lỗi định tuyến:} default route từ Site 500 làm mọi site mất Internet (thêm bộ lọc BGP), và Site 100 mất Internet do firewall không chia tải qua hai interface (điều chỉnh cost OSPF).
    \item \textbf{Campus NetOps Bot:} trợ lý vận hành qua Telegram, cài ngay trong topology và quản trị in-band Site 100 cùng 12 thiết bị chi nhánh qua overlay SD-WAN; chẩn đoán theo lớp, cảnh báo chủ động, thay đổi cấu hình có duyệt và tự rollback, hỏi đáp tự nhiên qua AI agent; selftest 23/23 ĐẠT.
\end{itemize}

Phần chưa làm theo kế hoạch: failover firewall HA, đo lặp SD-WAN đủ số lần và bản nháp báo cáo cuối kỳ. Tuần tới (02/11--08/11), nhóm tập trung vào các kịch bản sự cố còn lại (firewall, cắt một màu WAN), kiểm thử tích hợp Site 500, kịch bản demo NetOps Bot và bắt đầu viết báo cáo cuối kỳ.

\textbf{Từ khóa:} SD-WAN, Cisco Viptela, BFD, failover, SDN, Ryu, NetOps, ChatOps, Telegram, AI agent, EVE-NG.
